\documentclass[11pt, a4paper]{article}

\usepackage[style=script-v1, privacy=anonym]{../tools/library}

\renewcommand{\Lecturer}{Prof. S. Neukamm}
\begin{document}
\head{AN40: Differentialgleichungen \& Mannigfaltigkeiten}{Wintersemester 2025/2026}

\tableofcontents





\newpage
\section{Gewöhnliche Differentialgleichungen}

\vspace{0.3cm} \par \subsection{Grundlagen}


 \begin{definition} \Xlabel{Definition}{subsub}{Anfangswertprobleme (AWP)}

Ein \emph{Anfangswertproblem (AWP) $n$-ter Ordnung} auf einem Intervall $I\subset \R$ ist
\[
\quad \begin{cases}y^{(n)}(t)=f(t,y(t),\dots, y^{(n-1)}(t))\\ y^{(j)}(t_0)=y_j, \qquad 0\leq j\leq n-1\end{cases}
\]
wobei $f: I\times U_1 \times \dots \times U_{n-1}\to\R^d$ (mit $U_i\subset\R^d$ offen) gegeben ist. 
\end{definition}

\begin{note} \label{AWP} Ein AWP höherer Ordnung lässt sich durch Einführung von Zustandsvariablen auf ein äquivalentes System erster Ordnung reduzieren. Deshalb reicht es, Anfangswertprobleme 1. Ordnung der Form 

$$ (*)\quad \begin{cases}y'(t)=f(t,y(t))\\ y(t_0)=y_0 \end{cases} $$

zu studieren.
\end{note}
 \begin{example}
Reduktion: $y''+p(t)y'+q(t)y=r(t)$ mit $y_1=y,\; y_2=y'$ $\Rightarrow$ System
\[
\begin{pmatrix}y_1\\y_2\end{pmatrix}'=
\begin{pmatrix}0&1\\-q(t)&-p(t)\end{pmatrix}
\begin{pmatrix}y_1\\y_2\end{pmatrix}+
\begin{pmatrix}0\\r(t)\end{pmatrix}.
\]
\end{example}

 \begin{definition}
Eine Funktion $y:I\to\R^d$ heißt \emph{Lösung} des AWP, falls $y$  n-mal (stetig) differenzierbar ist und die Gleichung in $(*)$ sowie Anfangsbedingung erfüllt.
\end{definition}



\vspace{0.3cm} \par \subsection{Existenz und Eindeutigkeit}
 \begin{theorem}\Xlabel{Satz}{subsub}{Peano}
Sei $(*)$ wie in \ref{AWP}. Sei $f$ stetig auf einer Umgebung von $(t_0,y_0)$. Dann existiert mindestens eine lokale Lösung $y: J \to \R^d $ des AWP auf einem Intervall $J$ um $t_0$.
\end{theorem}

 \begin{note}
Peano garantiert lediglich Existenz — Eindeutigkeit braucht stärkere Annahmen. 
\end{note}


 \begin{theorem}\Xlabel{Satz}{subsub}{Picard--Lindelöf}
Sei $(*)$ ein AWP und $f(t,y)$ auf einem Rechteck $R=[t_0-h,t_0+h]\times B_r(y_0)$ stetig und lokal Lipschitz in $y$ (d.\,h. $\exists L$ mit $\|f(t,y)-f(t,z)\|\le L\|y-z\|$). Dann existiert eine eindeutige Lösung des AWP auf einem Intervall $I\subset R$; sie entsteht als Fixpunkt der Picard-Iteration
\[
y_{n+1}(t)=y_0+\int_{t_0}^t f(s,y_n(s))\,ds.
\]
\end{theorem}

 \begin{corollary}
Stetige Abhängigkeit der Lösung vom Anfangswert: Kleine Änderungen in $y_0$ führen (lokal) zu kleinen Änderungen der Lösung (abschätzbar durch Gronwall).
\end{corollary}



\vspace{0.3cm} \par \subsection{Maximale Lösungen und Fortsetzbarkeit}
 \begin{definition} \label{def:max-Forsetzbarkeit}
Eine Lösung $y:I_{\max}\to\R^d$ heißt \emph{maximal}, falls für alle Lösungen $\hat{y}: \hat{I} \to \R^d$ des AWP gilt: $\hat{I}\subset I_{\max}$ 
\end{definition}


 \begin{theorem}
Unter den Picard-Annahmen besitzt jedes AWP genau eine lokale Lösung; diese besitzt eine eindeutige maximale Fortsetzung. Entweder ist das maximale Intervall unbeschränkt, oder die Lösung läuft auf den Rand der Definitionsmenge zu (Blow-up).
\end{theorem}


\vspace{0.3cm} \par \subsection{Abschätzungen und Stabilität}

 \begin{lemma}\Xlabel{Lemma}{subsub}{von Gronwall}
Sei $I = [t_0,t_1)$ und seien $\alpha,\beta \in C(I;\R)$.
\begin{enumerate}
\item (Differentielle Form) Ist $y \in C^1(I;\R)$ und erfüllt
\[
y'(t) \le \alpha(t) + \beta(t)\,y(t)
\quad \text{für alle } t \in (t_0,t_1),
\]
so gilt für alle $t \in I$
\[
y(t) \le y(t_0)\,e^{\int_{t_0}^{t} \beta(s)\,ds}
+ \int_{t_0}^{t} \alpha(s)\,
e^{\int_{s}^{t} \beta(r)\,dr}\,ds .
\]

\item (Integrale Form)
Ist $y \in C(I;\R)$, $\beta \ge 0$ und gilt
\[
y(t) \le \alpha(t) + \int_{t_0}^{t} \beta(s)\,y(s)\,ds
\quad \text{für alle } t \in I,
\]
so folgt für alle $t \in I$
\[
y(t) \le \alpha(t)
+ \int_{t_0}^{t} \alpha(s)\,\beta(s)\,
e^{\int_{s}^{t} \beta(r)\,dr}\,ds .
\]
\end{enumerate}
\end{lemma}

 \begin{corollary}
Gronwall liefert die klassische Stetigkeitsabschätzung für Lösungen und die Lipschitz-Abhängigkeit von Anfangsbedingungen.
\end{corollary}



\begin{theorem}[Stetige Abhängigkeit vom Anfangswert]
Sei $f$ stetig und lokal Lipschitz-stetig bezüglich $y$. 
Sei $y:[t_0,t_1]\to\R^d$ eine Lösung zum Anfangswert $y_0$.
Dann existieren Konstanten $\delta,C>0$, sodass für alle
$\tilde y_0\in B(y_0;\delta)$ eine Lösung
$\tilde y:[t_0,t_1]\to\R^d$ existiert und
\[
\lvert y(t)-\tilde y(t)\rvert
\le C\,\lvert y_0-\tilde y_0\rvert
\quad \forall\, t\in[t_0,t_1].
\]
\end{theorem}


\begin{theorem}[Differenzierbare Abhängigkeit vom Anfangswert]
Sei $f$ zusätzlich stetig differenzierbar bezüglich $y$.
Dann ist die Lösungsabbildung
\[
p \longmapsto X(t;p)
\]
differenzierbar. Die Ableitung
\[
v(t)=\partial_p X(t;y_0)\,v_0
\]
erfüllt die linearisierte Differentialgleichung
\[
v'(t)=A(t)\,v(t), \qquad v(t_0)=v_0,
\]
wobei
\[
A(t):=\partial_y f\bigl(t,y(t)\bigr).
\]
\end{theorem}

\newpage
\section{Lineare Systeme}

\vspace{0.3cm} \par \subsection{Lineare Autonome Systeme}


\begin{definition}
    Betrachte wieder $(*)$ \ref{AWP}. Ein System der Form $f(t,y) = Ay(t)$ mit einer \underline{konstanten} Matrix $A \in \mathbb{R}^{d \times d}$
heißt lineares, homogenes, autonomes System.
\end{definition}


\begin{note} Linear-homogene Autonome Systeme sind so gutartig, dass wir ihre Lösungen direkt hinschreiben können. Insbesondere gelten die folgenden Eigenschaften.
    
\begin{enumerate}
\item \emph{Superpositionsprinzip}: Sind $y_1, y_2$ Lösungen, so ist auch $c_1 y_1 + c_2 y_2$ eine Lösung.
\item \emph{Lösungsraum}: Die Menge aller Lösungen bildet einen Vektorraum der Dimension $d$. Qualitativ formuliert: Sei $y'=Ay$ ein linear-homogenes autonomes System. Dann ist
$$\mathbb{L}=\left\{y_j = \exp(tA)e_j\; | \; 1\leq j \leq d = \sqrt{\dim(A)}\right\} $$
genau dieser d-dim. Lösungsraum. Für $\exp(tA)$ siehe \ref{Matrix-Exponential}
\item Globale Existenz \& Eindeutigkeit: Für jeden AW $y(t_0)=y_0$ existiert eine eindeutige, globale Lösung auf ganz $\mathbb{R}$.
\end{enumerate}
\end{note}


\begin{definition}\Xlabel{Definition}{subsub}{Matrix-Exponential}
Für eine Matrix $A \in \mathbb{C}^{d \times d}$ ist die Matrix-Exponentialfunktion definiert
durch die absolut konvergente Reihe:

\[
\exp(A) := \sum_{k=0}^{\infty} \frac{A^k}{k!}
= I + A + \frac{A^2}{2} + \dots
\]
\end{definition}


\begin{theorem} \label{lin-hom-Lsg}
Sei $A \in \mathbb{R}^{d \times d}$. Das AWP $y'=Ay$, $y(t_0)=y_0$
besitzt die eindeutige globale Lösung

\[
y(t) := \exp((t-t_0)A)y_0
\]
\end{theorem}


\begin{lemma} Sei $y(t)$ Lösung aus $\ref{lin-hom-Lsg}$. Dann gilt
\[
\frac{d}{dt}\exp(tA) = A\exp(tA) = \exp(tA)A
\]
\end{lemma}


\begin{example} Sei $y'(t)=Ay(t)$ DGL und $A$ ist reell diagonalisierbar. Dann gilt 
$$A = SDS^{-1}$$ für ein $S$ invertierbar und $D=\mathrm{diag}(\lambda_1,\dots,\lambda_d)$. Berechne nun $\exp(tA)$ mithilfe der Eigenschaft
$$\exp(SBS^{-1}) = S\exp(B)S^{-1}$$
\[
\Rightarrow\exp(tA) = \exp(tSDS^{-1}) = S\exp(tD)S^{-1}
\]

\[
= S
\begin{pmatrix}
e^{\lambda_1 t} & & \\
&\ddots& \\
& &e^{\lambda_d t}
\end{pmatrix}
S^{-1}
\]

Die Lösung ist eine Linearkombination von Termen der Form $e^{\lambda_j t} v_j$,
wobei $v_j$ die Eigenvektoren (Spalten von $S$) sind.
\end{example}

\begin{example} Sei die DGL gegeben durch den harmonischen Oszillator.
\[
y'' + y = 0 \;\Longleftrightarrow\; z' = Az
\]

mit $z=(y,y')$ und

\[
A=
\begin{pmatrix}
0 & 1 \\
-1 & 0
\end{pmatrix}
\]
A ist komplex diagonalisierbar.

$$A=\begin{pmatrix}
0 & 1 \\
-1 & 0
\end{pmatrix}
 =\begin{pmatrix}
-i & i \\
1 & 1
\end{pmatrix} \begin{pmatrix}
i & 0 \\
0 & -i
\end{pmatrix} \begin{pmatrix}
i/2 & 1/2 \\
-i/2 & 1/2
\end{pmatrix}$$
Für $z(0) = \vect{1\\0}$
\[
z(t) = \exp(tA)z(0)
=
\begin{pmatrix}
\cos(t) \\
-\sin(t)
\end{pmatrix}
\]

Reelle Lösungen entstehen aus komplex konjugierten Eigenwertpaaren
$\lambda=a\pm i\omega$ und haben die Form
$e^{at}(\cos(\omega t), \sin(\omega t))$.
\end{example} 

\begin{example} Sei $y'(t)=Ay(t)$ DGL und $A$ ist nicht diagonalisierbar. Dann ist dennoch jede Matrix $A$ ähnlich zu einer Matrix in Jordanscher Normalform
$J=S^{-1}AS$.

\[
J=\mathrm{diag}(J_1,\dots,J_\ell), \quad J_j=\lambda_j I + N_j
\]

wobei $N_j$ ist eine nilpotent ist. Nehme
\[
A=
\begin{pmatrix}
0 & 1 & 0 \\
0 & 0 & 1 \\
0 & 0 & 0
\end{pmatrix}
\]

$A$ ist nilpotent ($A^3=0$).

\[
\exp(tA) = I + tA + \frac{t^2}{2}A^2
=
\begin{pmatrix}
1 & t & t^2/2 \\
0 & 1 & t \\
0 & 0 & 1
\end{pmatrix}
\]

Allgemein:
\[
\exp(tJ_j)=e^{\lambda_j t}\exp(tN_j)
\]

Lösungen enthalten also Terme der Form $t^k e^{\lambda t}$.
\end{example}



\vspace{0.3cm} \par \subsection{Lineare Zeitvariante Systeme}

\begin{definition}
Ein System der Form $y'=A(t)y+b(t)$ mit stetigen, zeitabhängigen
$A(t), b(t): \R^d \to \R^d$ heißt lineares, \underline{zeitvariantes} System.
\end{definition}

\begin{theorem}\Xlabel{Satz}{subsub}{Globale Existenz und Eindeutigkeit}
Sei $I\subset\mathbb{R}$ ein offenes Intervall, $A(t)$ und $b(t)$ stetig auf $I$.
Dann hat das AWP
$y'=A(t)y+b(t)$, $y(t_0)=y_0$
für jedes $(t_0,y_0)\in I\times\mathbb{R}^d$
eine eindeutige, globale Lösung $y:I\to\mathbb{R}^d$. (folgt aus \ref{Picard--Lindelöf} und \ref{von Gronwall})
\end{theorem}



\begin{theorem}\Xlabel{Satz \& Definition}{subsub}{Fundamentalmatrix} Sei $A:I\to\mathbb{R}^{d\times d}$ stetig auf $I$ offen.
Es existiert genau eine Abbildung
\[
\Phi:I\times I \to \mathbb{R}^{d\times d}
\]

sodass für alle $t_0\in I$, $y_0\in\mathbb{R}^d$ die Abbildung

\[
y(t):=\Phi(t,t_0)y_0
\]
das homogene AWP $y'=A(t)y$, $y(t_0)=y_0$ löst.
$\Phi$ heißt \underline{Fundamentalmatrix} und hat folgende Eigenschaften:

\begin{enumerate}
\item $\Phi(t,t)=I$
\item Funktionalgleichung: $\Phi(t,s)=\Phi(t,r)\Phi(r,s)$
\item Ableitung: $\frac{d}{dt}\Phi(t,s)=A(t)\Phi(t,s)$
\item Invertierbarkeit: $\Phi(t,s)^{-1}=\Phi(s,t)$
\end{enumerate}

Falls $A$ nicht von $t$ abhängt:
\[
\Phi(t,s)=\exp((t-s)A)
\]

Falls $A(t)=a(t)\in\mathbb{R}$:
\[
\Phi(t,s)=\exp\left(\int_s^t a(\tau)d\tau\right)
\]
    
\end{theorem}



\begin{lemma}\Xlabel{Lemma}{subsub}{Lineare Inhomogene Systeme} Mihilfe von Variation der Konstanten. \bigskip

Sei $\Phi(t,s)$ die Fundamentalmatrix von $y'=A(t)y$.
Die Lösung des inhomogenen AWP
$y'=A(t)y+b(t)$, $y(t_0)=y_0$
ist gegeben durch:

\[
y(t)=\Phi(t,t_0)y_0
+
\int_{t_0}^{t} \Phi(t,s)b(s)\,ds
\]

\end{lemma}

\begin{lemma}\Xlabel{Lemma}{subsub}{Skalare homogene Systeme 1. Ordnung}

Betrachte das skalare AWP $y'=f(t)g(y)$, $y(t_0)=y_0$.

\begin{enumerate}
\item Fall: Wenn $g(y_0)=0$, ist $y(t)\equiv y_0$ eine konstante Lösung.
\item Fall: Wenn $g(y_0)\neq 0$, so ist die Lösung implizit gegeben durch:
\end{enumerate}

\[
\int_{y_0}^{y(t)} \frac{1}{g(x)}dx
=
\int_{t_0}^{t} f(s)ds
\]
\end{lemma}


\newpage
\section{Untermannigfaltigkeiten}

\vspace{0.3cm} \par \subsection{Grundlagen}


\begin{definition}\Xlabel{Definition}{subsub}{Parametrisierte Kurve} Eine stetige Abb. $\gamma : I \to \mathbb{R}^d$ heißt parametrisierte Kurve. $\gamma$ heißt regulär, wenn
$\gamma \in C^1$ und $\gamma'(t) \neq 0$.
\end{definition}


\begin{definition}\Xlabel{Definition}{subsub}{Länge einer \texorpdfstring{$C^1$}{C¹}-Kurve}. Sei $\gamma\in C^1(I,\R^d)$. Dann bezeichnet
\[
L(\gamma) = \int_a^b |\gamma'(t)|_2 \, dt.
\]
die Länge von $\gamma$.
\end{definition}

\begin{definition}\Xlabel{Definition}{subsub}{Kurvenintegral}. Für ein Vektorfeld $F$ ist
\[
\int_\gamma F \cdot dx
:=
\int_a^b F(\gamma(t)) \cdot \gamma'(t) \, dt.
\]
Das sogenannte Kurvenintergral unter $\gamma$
\end{definition}


\begin{example} Sei 
\[
\gamma(t) = (r\cos t, r\sin t, ct), \quad t \in [0,2\pi].
\]
$\gamma$ ist $ C^1$.
\[
\gamma'(t) = (-r\sin t, r\cos t, c)
\quad \Rightarrow \quad
|\gamma'(t)|_2 = \sqrt{r^2 + c^2}.
\]

Es ergeben sich Länge daher die Länge über dem Intervall $[0,2\pi]$:

\[
L(\gamma)
=
\int_0^{2\pi} \sqrt{r^2 + c^2}\, dt
=
2\pi \sqrt{r^2 + c^2}.
\]

Sei $F=(0,0,-1)$ ein Vektorfeld. Dann ist die physikalisch verrichtete Arbeit entlang von $\gamma$:

\[
\int_\gamma F \cdot dx
=
\int_0^{2\pi} (0,0,-1)\cdot \gamma'(t)\, dt
=
-2\pi c.
\]

\end{example}

\begin{theorem} \Xlabel{Satz}{subsub}{Wegunabhängigkeit und Gradientenfelder} Sei $\Omega \subset \mathbb{R}^d$ offen und wegzusammenhängend,
$F:\Omega \to \mathbb{R}^d$ stetig.

Äquivalent sind:

\begin{itemize}
\item[(a)] $F$ ist ein Gradientenfeld: $\exists f \in C^1(\Omega)$ mit $F=\nabla f$.
\item[(b)] $F$ ist konservativ (wegunabhängig):
\[
\int_{\gamma_1} F\cdot dx
=
\int_{\gamma_2} F\cdot dx
\]
für alle Kurven $\gamma_1,\gamma_2$ mit gleichen Endpunkten.
\end{itemize}
\end{theorem}


\begin{lemma} Sei $F=\nabla f$, dann gilt:

\[
\int_\gamma \nabla f \cdot dx
=
f(\gamma(b)) - f(\gamma(a)).
\]

\end{lemma}




\vspace{0.3cm} \par \subsection{Konstruktion}

\begin{definition} \Xlabel{Definition}{subsub}{Immersion}
Seien $U \subset \mathbb{R}^k$ offen und $\varphi : U \to \mathbb{R}^n$ eine $C^1$-Abbildung. Die Abbildung $\varphi$ heißt \emph{Immersion}, wenn die Jacobi-Matrix $Df(x)$ für alle $x \in U$  vollen Rang $k$ besitzt, d.h.
\[
\operatorname{rang}(Df(x)) = k 
\quad \text{für alle } x \in U.
\]
    
\end{definition}

\begin{definition} \Xlabel{Definition}{subsub}{Einbettung}
    Sei $\varphi \in C^1(U;\R^n)$ eine Immersion. Falls 
    $$\psi: U \to \varphi(U),\quad x \mapsto \varphi(x)$$ ein Homöomorphismus ($U\cong \varphi(U))$ ist, dann heißt $\psi$ \emph{Einbettung}.
\end{definition}

\begin{note} $\psi$ Einbettung ist ein lokaler Diffeomorphismus, denn $$D\psi = D\varphi$$
\end{note}


\begin{definition} \Xlabel{Definition}{subsub}{Untermannigfaltigkeit} $M \subset \mathbb{R}^n$ ist eine $k$-dimensionale $C^1$-Untermannigfaltigkeit,
wenn für jedes $a\in M$ lokal eine der äquivalenten Darstellungen gilt:

\begin{itemize}
\item[(a)] Parameterdarstellung: $M$ ist lokal Bild einer Einbettung
$\psi:U \to \mathbb{R}^n$.

\item[(b)] Implizite Darstellung: $M$ ist lokal Nullstellenmenge einer Funktion\\  $f: \R^n\supset \widehat V \to \mathbb{R}^{n-k}$ mit vollem Rang $(\rang{Df} = n-k \; \forall x \in \widehat V)$

\item[(c)] Kartendarstellung: $M$ ist lokal diffeomorph zu $\mathbb{R}^k$.
\end{itemize}

Weitere Darstellung: $M$ ist lokal Graph einer $C^1$-Funktion.
\end{definition}

\begin{definition}\Xlabel{Definition}{subsub}{Karte, Atlas} Sei $M\subset \R^n$ eine $k-$ dim. $C^1-$ Untermannigfaltigkeit. Wir nennen $V\subset M$ offen ein \underline{Kartengebiet}, falls es $\widehat V , \widehat U$ und ein $\widehat \Phi: \widehat V \to \widehat U$ Diffeomorphismus gibt, sodass $\widehat V \cap M = V$ und $\widehat U \cap \R^n_k = \widehat \Phi (\widehat V)$. In diesem Fall heißt $\widehat \Phi$ \underline{Karte} für $M$. Eine Familie von Karten $(\Phi_i)_{i\in I}$ zu Kartengebieten $(V_i)_{i\in I}$ heißt \underline{Atlas} von $M$, falls $\bigcup_{i\in I} V_i \supset M$.
\end{definition}
\begin{example} \Xlabel{Beispiel}{subsub}{Sphäre} Betrachte 
\[
S^{n-1}
=
\{x \in \mathbb{R}^n : |x|_2^2 = 1\}
\]
ist eine $(n-1)$-dimensionale Untermannigfaltigkeit.\bigskip

Implizite Darstellung:

\[
f(x)=|x|_2^2 -1 \quad \& \quad \nabla f(x)=2x \neq 0 \quad \text{für } x\in S^{n-1}.
\]

Rang der Jacobi-Matrix $Jf(x)$ ist
\[
1 = n-(n-1).
\]

$S^{n-1}$ ist Nullstellenmenge von $f$.

\end{example}


\vspace{0.3cm} \par \subsection{k-dimensionales Volumen}



\begin{definition} \Xlabel{Definition}{subsub}{Maßtensor} Sei $\varphi: U \to \R^k$ eine Immersion. Die Abbildung
$$G_\varphi: U \to \R^{k\times k}, \quad y\mapsto J_\varphi(y)^\top J_\varphi (y)$$
heißt \underline{Maßtensor} von $\varphi$ und 
$$g^\varphi := \det G_\varphi$$
die \underline{Gram'sche Determinante} von $\varphi$.
\end{definition}

\begin{definition} \Xlabel{Definition}{subsub}{k-dim. Volumina}
    Sei $\varphi: U \to \R^n$ eine Immersion und $V\subset U$ meßbar. Dann heißt
    $$\Vol_k(\varphi|_V) := \int _V \sqrt{g^\varphi(y)} \lambda^k(dy)$$
    das \underline{$k-$dim. Volumen} von $\varphi$.\bigskip

    Falls $\exists \varphi: U \to \R^n$ Immersion, wobei $\varphi(U)\cong W$ homöomorph, so nennen wir 
    $$\Vol_k(W) = \Vol_k(\varphi) $$
    das  \underline{$k-$dim. Volumen} von $W$.
\end{definition}

\begin{note} Falls $\varphi_1, \varphi_2$ Einbettungen sind, die $W$ parametrisieren, dann gilt:
$$\Vol_k(\varphi_1) = \Vol_k(W) = \Vol_k(\varphi_2)$$
\end{note}

\vspace{0.3cm} \par \subsection{Oberflächenintegral}
\begin{definition}\Xlabel{Definition}{subsub}{Oberflächenintegral über ein Kartengebiet} Sei $V$ ein Kartengebiet einer $k-$ dim. Untermannigfaltigkeit der $\R^n$. Und $u:V\to \RT$ eine meßbare Funktion. Dann heißt $u$ integrierbar auf $V$, falls für eine Einbettung $\varphi: U \to \R^n$ mit $U\subset \R^k , \varphi(U) = V$ die Funktion $(u\circ \varphi)\sqrt{g^\varphi}$ (Lebesgue-)integrierbar auf $U$ ist. In diesem Fall heißt:

$$\int_V u \d S := \int_U u(\varphi(y))\sqrt{g^\varphi(y)} dy$$
Das Oberflächenintegral von $u$ auf $V$.
\end{definition}

\begin{example}
    Sei $V$ ein Kartengebiet einer $k$-dim. Untermannigfaltigkeit $M$ und $\varphi: U \to \R^n$ eine Einbettung mit $V= \varphi(U)$.

    \begin{enumerate}
        \item Sei $W\subset V$ offen in $V$. Dann ist $U':=\varphi\r (W)$ offen in $U$ ($\varphi$ homöomorph) und damit $U\ni y\mapsto \1 _{U'}(y)\sqrt{g^\varphi (y)}$ meßbar. Es gilt
        $$\int _V \1 _{W}(x)\d S(x) = \int _U \1 _{W}(\varphi(y))\sqrt{g^\varphi (y)} dy= \int _U \1 _{U'}(y)\sqrt{g^\varphi (y)} dy $$
        $$= \int _{U'} \sqrt{g^\varphi (y)} dy= \Vol_k(\varphi|_{U'})$$
        Da aber $\varphi|_{U'}: U' \to \R^n$ eine Einbettung mit $\varphi(U')= W$ ist, gilt nach \ref{k-dim. Volumina}, dass
        $$\Vol_k(W) = \int_V \1_W \d S$$
        \item Sei $V = \Gamma(f)$ Graph einer Funktion $f \in C^1(U)$. Sei $u : V \to \R $ integrierbar. Dann gilt
$$ \int _V u \d S = \int _U u(y,f(y)) \sqrt{1+ |\nabla f(y)|^2 _2 }\d y$$
    \end{enumerate}
    
\end{example}


\begin{definition}\Xlabel{Definition}{subsub}{Zerlegung der Eins}
Sei $M$ eine Untermannigfaltigkeit. Unter einer \emph{Zerlegung der Eins} auf $M$ verstehen wir eine Familie stetiger Funktionen $(\eta_i)_{i \in I} \subset C(M)$ mit abzählbarer Indexmenge $I$ und folgenden Eigenschaften:

\begin{enumerate}
    \item[(i)] $0 \le \eta_i(x) \le 1$ für alle Punkte $x \in M$ und alle $i \in I$.
    
    \item[(ii)] Die Familie $(\eta_i)_{i \in I}$ ist lokal endlich, d.\,h. für jeden Punkt $x \in M$ existiert eine (in $M$) offene Umgebung $V$ von $x$, sodass $\eta_i = 0$ in $V$ für alle bis auf endlich viele $i \in I$.
    
    \item[(iii)] Es gilt 
    \[
        \sum_{i \in I} \eta_i(x) = 1
    \]
    für alle Punkte $x \in M$.
\end{enumerate}

Ist $(V_i)_{i \in I}$ eine offene Überdeckung von $M$ (z.\,B. ein abzählbarer Atlas von $M$), so sagen wir, dass die Zerlegung der Eins $(\eta_i)_{i \in I}$ der Überdeckung $(V_i)_{i \in I}$ \emph{untergeordnet} ist, falls zusätzlich gilt:

\begin{enumerate}
    \item[(iv)] Für jedes $i \in I$ ist der Träger $\operatorname{supp}(\eta_i)$ in $V_i$ enthalten.
\end{enumerate}
\end{definition}

\begin{theorem} Jede Untermannigfaltigkeit des $\R^n$ besitzt einen
abzählbaren Atlas $(V_i,\varphi_i)_{i\in I}$ und eine untergeordnete Zerlegung der Eins ${\eta _i }_{i\in I}$
\end{theorem}

\begin{definition}\Xlabel{Definition}{subsub}{Oberflächenintegral einer UMf. $M$}
Sei $M \subset \mathbb{R}^n$ eine Untermannigfaltigkeit. Eine Funktion 
$u : M \to \overline{\mathbb{R}}$ heißt \emph{messbar auf $M$}, wenn ein abzählbarer Atlas 
$(V_i)_{i \in I}$ und eine diesem untergeordnete Zerlegung der Eins 
$(\eta_i)_{i \in I}$ existieren, so dass die Funktionen $u \eta_i$, $i \in I$, 
messbar auf $M$ sind.

Gilt zusätzlich
\[
\sum_{i \in I} \int_{V_i} |u| \eta_i \, dS < \infty,
\]
so heißt $u$ \emph{integrierbar auf $M$}, und
\[
\int_M u \, dS 
:= 
\sum_{i \in I} \int_{V_i} u \eta_i \, dS
\]
heißt das \emph{Oberflächenintegral} von $u$ auf $M$.
\end{definition}
    






\vspace{0.3cm} \par \subsection{Tangentialräume}
\begin{definition}
\Xlabel{Definition}{subsub}{Tangentialvektor und Tangentialraum} Sei $M \subset \mathbb{R}^n$ eine Untermannigfaltigkeit und $a \in M$. Ein Vektor $v \in \mathbb{R}^n$ heißt \emph{Tangentialvektor an $M$ in $a$}, 
falls eine $C^1$-Kurve $\gamma : (-\delta,\delta) \to M$ mit $\delta > 0$ existiert, so dass
\[
\gamma(0) = a 
\quad \text{und} \quad 
\gamma'(0) = v.
\]

Die Menge aller Tangentialvektoren an $M$ in $a$ heißt der 
\emph{Tangentialraum von $M$ in $a$} und wird mit
\[
T_a (M)
\]
bezeichnet.
\end{definition}

\begin{theorem}
\Xlabel{Satz}{subsub}{Dimension des Tangentialraums}

Sei $M \subset \mathbb{R}^n$ eine $k$-dimensionale Untermannigfaltigkeit 
und $a \in M$. Dann ist der Tangentialraum $T_a M$ ein 
$k$-dimensionaler Untervektorraum von $\mathbb{R}^n$.
\end{theorem}

\begin{definition}
\Xlabel{Definition}{subsub}{Normalenraum}

Sei $M \subset \mathbb{R}^n$ eine $(n-1)$-dimensionale Untermannigfaltigkeit 
und $a \in M$. Der \emph{Normalenraum} von $M$ in $a$ ist das orthogonale Komplement des Tangentialraums:
\[
N_a M := (T_a M)^\perp.
\]
\end{definition}




\begin{example}

Sei $M \subset \mathbb{R}^n$ implizit durch eine $C^1$-Abbildung 
$f : \mathbb{R}^n \to \mathbb{R}^m$ gegeben mit
\[
M = \{x \in \mathbb{R}^n : f(x) = 0\},
\]
und sei $a \in M$. Dann gilt
\[
T_a M 
=
\mathrm{span}\{\nabla f_1(a), \dots, \nabla f_m(a)\}^\perp.
\]

Insbesondere für die Einheitssphäre
\[
S^{n-1}
=
\{x \in \mathbb{R}^n : f(x) = 0\}
\quad \text{mit} \quad
f(x) = |x|^2 - 1
\]
ist
\[
\nabla f(x) = 2x.
\]

Für $a \in S^{n-1}$ folgt daher
\[
T_a S^{n-1}
=
\{ v \in \mathbb{R}^n : a \cdot v = 0 \}.
\]

Der Tangentialraum ist also die Hyperebene durch den Ursprung, 
die orthogonal auf dem Vektor $a$ steht.
\end{example}

\begin{definition}
\Xlabel{Definition}{subsub}{Einheitsnormalenfeld} Sei $M \subset \mathbb{R}^n$ eine $(n-1)$-dimensionale Untermannigfaltigkeit. Unter einem \emph{Einheitsnormalenfeld} auf $M$ verstehen wir eine \underline{stetige} Abbildung
\[
\nu : M \to \mathbb{R}^n,
\]
so dass für alle $x \in M$ gilt:
\[
\nu(x) \perp T_x M
\quad \text{und} \quad
|\nu(x)|^2 = 1.
\]
\end{definition}

\begin{lemma} Sei $M \subset \mathbb{R}^n$ eine $(n-1)$-dimensionale, wegzusammenhängende 
Untermannigfaltigkeit. Dann gibt es auf $M$ entweder kein 
oder genau zwei Einheitsnormalenfelder.
\end{lemma}

\begin{definition}
\Xlabel{Definition}{subsub}{Orientierbarkeit} Eine $(n-1)$-dimensionale Untermannigfaltigkeit $M \subset \mathbb{R}^n$ heißt \emph{orientierbar}, 
falls sie ein Einheitsnormalenfeld $\nu$ besitzt. In diesem Fall nennt man $(M,\nu)$ eine 
\emph{orientierte $(n-1)$-dimensionale Untermannigfaltigkeit} 
im $\mathbb{R}^n$.
\end{definition}

\begin{example}

\begin{enumerate}
    \item[(a)] Die Sphäre
    \[
    S^{n-1} := \{ x \in \mathbb{R}^n : |x|^2 = 1 \}
    \]
    ist orientierbar. Die beiden Einheitsnormalenfelder sind gegeben durch
    \[
    \nu_\pm(x) := \pm x.
    \]

    \item[(b)] Das Möbiusband ist nicht orientierbar .
\end{enumerate}
\end{example}


\vspace{0.3cm} \par \subsection{Globale Integralsätze}

\begin{definition}
\Xlabel{Definition}{subsub}{$C^1$-Rand} Sei $\Omega \subset \mathbb{R}^n$ offen. 
Wir sagen, dass $\Omega$ einen \emph{$C^1$-Rand} besitzt, 
falls für alle $a \in \partial \Omega$ eine offene Menge 
$U \subset \mathbb{R}^n$ und eine Funktion $f \in C^1(U)$ 
existieren, so dass $\nabla f(x) \neq 0$ für alle $x \in U$ gilt und

\begin{enumerate}[label=\alph*)]
\item $\partial \Omega \cap U = \{ x \in U : f(x) = 0 \}$
\item $\Omega \cap U = \{ x \in U : f(x) < 0 \}$
\end{enumerate}

Ist $\Omega \subset \mathbb{R}^n$ ein Gebiet mit $C^1$-Rand, 
so liegt $\Omega$ lokal um jeden Randpunkt $x \in \partial \Omega$ 
immer nur auf einer Seite des Randes. 

Außerdem ist der Rand so regulär (nämlich eine $C^1$-Untermannigfaltigkeit), 
dass an jedem Randpunkt eine eindeutig bestimmte Normale existiert.
\end{definition}

\begin{lemma} Sei $\Omega \subset \mathbb{R}^n$ offen mit $C^1$-Rand.  
Dann ist $\partial \Omega$ orientierbar und es existiert genau ein 
Einheitsnormalenfeld 
\[
\nu : \partial \Omega \to \mathbb{R}^n
\] 
an $\partial \Omega$, das nach außen zeigt, d.h.
\[
\forall a \in \partial \Omega \;\; \exists \varepsilon > 0 : 
a + t \, \nu(a) \notin \Omega \quad \text{für alle } t \in (0, \varepsilon).
\]

Der Vektor $\nu(a)$ heißt \emph{äußerer (Einheits-)Normalvektor} an 
$\partial \Omega$ in $a$.
\end{lemma}

\begin{note} Sei $\Omega \subset \mathbb{R}^n$ offen mit $C^1$-Rand und lokal beschrieben durch
\[
\Omega \cap U = \{ x \in U : f(x) < 0 \}, \qquad 
\partial \Omega \cap U = \{ x \in U : f(x) = 0 \}, 
\]
mit $\nabla f(a) \neq 0$ für alle $a \in \partial \Omega$.  

Dann ist der äußere Einheitsnormalenvektor $\nu(a)$ genau gegeben durch den normierten Gradienten:
\[
\nu(a) = \frac{\nabla f(a)}{|\nabla f(a)|_2}.
\]

Dieses Vorzeichen des Gradienten wählt die Richtung „nach außen“, also aus $\Omega$ heraus.
\end{note}


\begin{note}
\Xlabel{Definition}{subsub}{Divergenz und Rotation eines Vf.}

Sei $F = (F_1, \dots, F_n)$ ein $C^1$-Vektorfeld auf einem offenen Gebiet in $\mathbb{R}^n$. Dann ist definiert

\begin{enumerate}[label = $-$]
    \item \emph{Divergenz:} 
    \[
    \operatorname{div} F := \sum_{i=1}^n \frac{\partial F_i}{\partial x_i}.
    \]
    Sie misst die lokale Quell- bzw. Senkenstärke des Feldes.

    \item \emph{Rotation (für $n=3$):} 
    \[
    \operatorname{rot} F := 
    \nabla \times F 
    =
    \begin{pmatrix}
    \frac{\partial F_3}{\partial x_2} - \frac{\partial F_2}{\partial x_3} \\
    \frac{\partial F_1}{\partial x_3} - \frac{\partial F_3}{\partial x_1} \\
    \frac{\partial F_2}{\partial x_1} - \frac{\partial F_1}{\partial x_2}
    \end{pmatrix}.
    \]
    Sie misst die lokale Wirbelstärke des Feldes.
\end{enumerate}

Für $n \neq 3$ kann die Rotation in \(\mathbb{R}^n\) über Differentialformen verallgemeinert werden,  
in \(\mathbb{R}^2\) reduziert sich \(\operatorname{rot} F\) auf einen Skalar: 
\(\partial_1 F_2 - \partial_2 F_1\).
\end{note}


\begin{theorem}
\Xlabel{Satz}{subsub}{Integralsatz von Gauß}

Sei $\Omega \subset \mathbb{R}^n$ offen und beschränkt, und sei $\Omega$ ein Gebiet mit $C^1$-Rand.  
Sei $U \subset \mathbb{R}^n$ offen mit $\overline{\Omega} \subset U$ und sei 
$F : U \to \mathbb{R}^n$ ein $C^1$-Vektorfeld.  

Dann gilt:
\[
\int_{\Omega} \operatorname{div} F(x) \, dx
=
\int_{\partial \Omega} F(a) \cdot \nu(a) \, dS(a),
\]
wobei $\nu : \partial \Omega \to \mathbb{R}^n$ das äußere Einheitsnormalenfeld bezeichnet.
\end{theorem}

\begin{note}

Im Spezialfall $n=1$, bei zusammenhängendem $\Omega$, ist $\Omega$ gerade ein offenes Intervall 
der Form $\Omega = (a,b)$, $F$ ist ein Skalarfeld und $\operatorname{div} F = F'$.  

In diesem Fall liefert der Satz von Gauß gerade den Hauptsatz der Integral- und Differentialrechnung:
\[
\int_a^b F'(x) \, dx = F(b) - F(a).
\]

Wir können uns daher den Satz von Gauß als eine Verallgemeinerung 
des Hauptsatzes der Differential- und Integralrechnung auf höhere Dimensionen vorstellen.
\end{note}


\begin{lemma}
\Xlabel{Lemma}{subsub}{1. und 2. Greensche Formel} Sei $\Omega \subset \mathbb{R}^n$ offen und beschränkt mit $C^1$-Rand.  
Sei $f$ (bzw. $g$) eine in einer offenen Umgebung von $\overline{\Omega}$ einmal (bzw. zweimal) stetig differenzierbare Funktion. Dann gilt die \emph{erste Greensche Formel}:
\[
\int_{\Omega} f(x) \, \Delta g(x) \, dx
=
- \int_{\Omega} \nabla f(x) \cdot \nabla g(x) \, dx
+ \int_{\partial \Omega} f(a) \, (\nabla g(a) \cdot \nu(a)) \, dS(a).
\]

Im Spezialfall $f = 1$ folgt:
\[
\int_{\Omega} \Delta g(x) \, dx
=
\int_{\partial \Omega} \nabla g(a) \cdot \nu(a) \, dS(a).
\]

Ist $f$ zusätzlich zweimal stetig differenzierbar in einer offenen Umgebung von $\overline{\Omega}$, 
dann gilt die \emph{zweite Greensche Formel}:
\[
\int_{\Omega} \bigl( f(x) \, \Delta g(x) - g(x) \, \Delta f(x) \bigr) dx
=
\int_{\partial \Omega} \bigl( f \, (\nabla g \cdot \nu) - g \, (\nabla f \cdot \nu) \bigr) dS.
\]

\end{lemma}


\begin{note}

Hierbei ist $\nabla g \cdot \nu$ gerade die Richtungsableitung von $g$ in Richtung der äußeren Normalen.  
Häufig verwendet man hierfür die Notation
\[
\partial_\nu g := \nabla g \cdot \nu.
\]
\end{note}


\begin{theorem}
\Xlabel{Satz}{subsub}{Satz von Stokes im $R^3$}

Sei $M \subset \hat{V}$ eine 2-dimensionale orientierbare Untermannigfaltigkeit in $\mathbb{R}^3$,  
die durch eine Einbettung 
\[
\varphi \in C^2(U ; \mathbb{R}^3), \quad U \subset \mathbb{R}^2 \text{ offen},
\] 
global parametrisiert wird.  

Sei $A \subset M$ eine beschränkte, offene reguläre Fläche mit Rand $\partial A$,  
der durch eine $C^1$-Kurve parametrisiert werden kann. Genauer:  
Es sei 
\[
\Omega := \varphi^{-1}(A) \subset U
\] 
offen und beschränkt, es gelte $\Omega$ beschränkt mit $C^1$-Rand,  
und es existiere eine reguläre Kurve
\[
\gamma : [0,1] \to \mathbb{R}^2
\] 
mit $\gamma(0) = \gamma(1)$, die das Intervall $[0,1)$ bijektiv auf $\partial \Omega$ abbildet.  

Wir bezeichnen das äußere Einheitsnormalenfeld mit
\[
\nu : \partial \Omega \to \mathbb{R}^2.
\] und $\overline{\Omega} \subset U$.  

Dann gilt für jedes $C^1$-Vektorfeld 
\[
F : \hat{U} \to \mathbb{R}^3
\] 
auf einer offenen Umgebung $\hat{V} \subset \mathbb{R}^3$ von $M$ die Identität:
\[
\int_{\partial A} F(x) \cdot dx
=
\int_A \operatorname{rot} F(x) \cdot \nu(x) \, dS(x),
\]
wobei der Umlaufsinn der Randkurve des Kurvenintegrals auf der linken Seite  
und die Normalenrichtung $\nu(x)$ (wie in (10.5) und (10.7) beschrieben) zueinander passen. \bigskip

Wir schreiben in Kurzform auch mit dem Kurvenintegral $\oint$: \bigskip

In $\mathbb{R}^2$:

\[
\oint_{\partial\Omega} F\cdot dx
=
\int_\Omega
\left(
\frac{\partial F_2}{\partial x_1}
-
\frac{\partial F_1}{\partial x_2}
\right)
dx_1 dx_2.
\]

In $\mathbb{R}^3$:

\[
\oint_{\partial A} F\cdot dx
=
\int_A (\mathrm{rot}\,F)\cdot \nu\, dS.
\]
\end{theorem}

\PropertyNotice
\end{document}